\documentclass[10pt,a4paper]{amsart}
\usepackage[margin=19mm]{geometry}
\usepackage{amsmath,amssymb,mathtools}
\usepackage[scr=rsfs,cal=euler]{mathalfa}
\usepackage{xfrac}
\usepackage[unicode,hidelinks,bookmarks=false]{hyperref}
\newcommand{\ie}{i.e.\ }
\newcommand{\cf}{cf.\ }
\newcommand{\resp}{respectively\ }
\newcommand{\Subsection}{\subsection}
\providecommand{\nobreakdash}{\nobreak\hbox{-}}
\setlength{\emergencystretch}{2em}
\multlinegap0pt
\usepackage{amssymb}
\usepackage{float}
\usepackage{caption}
\usepackage{tikz}
\newtheorem{theorem}{Theorem}
\title{Propagation processes on (hyper)graphs:\\ where zero forcing and burning meet}
\author{Aida Abiad, Pax Mallee}
\date{Source: arXiv:2602.13149v1 (CC BY 4.0)}
\begin{document}
\maketitle

\noindent\emph{Source and licence.} Propagation processes on (hyper)graphs:\\ where zero forcing and burning meet. Version arXiv:2602.13149v1 (CC BY 4.0), distributed by arXiv under the Creative Commons Attribution 4.0 International licence, http://creativecommons.org/licenses/by/4.0/, as recorded in the arXiv licence metadata for this exact version and shown on the abstract page https://arxiv.org/abs/2602.13149v1. Full licence notice: \texttt{licenses/arxiv-2602.13149-CC-BY-4.0.txt}.\par\smallskip

\noindent\textit{Editorial scope.} Counted unit: the article's theorem inequality (source0.tex lines 976--983). The article's complete original proof of it is delivered with it (source0.tex lines 984--1012, one \texttt{proof} environment of 29 lines). No mathematics is written, repaired or re-derived: the statement and the proof are the article's own lines, in the article's own notation, and no retained line is substituted or re-flowed.  Retained uncounted as the source-local context the proof needs: lines 839--841.  Nothing was omitted from any retained span, so the package carries no omission record.  No retained line contains a citation, so the wrapper carries no bibliography and no bibliography entry.  No retained line contains a figure environment, an image-inclusion command or drawing code, so no image is delivered and the package declares no asset dependency.  Editorial formatting only, in addition: an AMS \texttt{amsart} preamble that restates the article's own declarations and macros used by the retained lines, namely the environments \texttt{theorem} on separate counters, together with the standard packages those lines need.  The retained environments print in this document as Theorem 1 in order of appearance, whereas the article prints the same items with the numbering of its own full document.  The complete native source of the article as distributed in the arXiv e-print for this exact version is bundled unchanged as \texttt{sources/194619/source0.tex}.
   \begin{equation}\label{eq:hypergraph inequality}
           Z(IG(H)) \leq b_L(H) + |E(H)|.
   \end{equation}

\begin{theorem}
    Let $H$ be a hypergraph with $k$ connected components that contain at least one non-singleton hyperedge. Then
    \begin{equation}
        Z(IG(H)) \leq b_L(H) + |E(H)| - k.
        \label{eq:component hypergraph inequality}
    \end{equation}
    \label{thm:improved bound}
\end{theorem}

\begin{proof}
    First, we show that $Z(IG(H_i)) \leq b_L(H_i) + |E(H_i)| - 1$ for each of the $k$ connected components $\{H_i\}_{1\leq i\leq k}$ containing at least one non-singleton hyperedge, which is done with proof by contradiction. Then, the relation is shown for the whole hypergraph.
    So, let $\{H_i\}_{1\leq i\leq l}$ be the set of all the $l\geq k$ connected components of $H$ such that $H_i$ for $1\leq i \leq k$ contains at least one non-singleton hyperedge and that $H_i$ for $k< i\leq l$ does contain none or only singleton hyperedges. 

    
\begin{description}
    \item[Case that $B_i\neq \emptyset$.] Take any $v\in B_i$ and take $h\in E(H_i)$ such that $v\in h$ and $h$ is non-singleton. Such $v$ and $h$ exist since $H_i$ is connected and contains at least one non-singleton hyperedge. 
    Earlier, it was shown that $B_i\cup E(H_i)$ was a zero forcing set for $IG(H_i)$. 
    Now, consider the set $B_i\cup E(H_i) \backslash \{h\}$ and we claim that it is a zero forcing set in $IG(H_i)$. 
    Since all the vertices of $IG(H_i)$ that are in $E(H_i)$ apart from $h$ are colored, and $v$ is connected in $IG(H_i)$ to vertices in $E(H_i)$, in the first round the force $v\rightarrow h$ will occur as $h$ is the only uncolored neighbor of $v$. 
    Then at least the set $B_i\cup E(H_i)$ is colored after the first round, which means that the set $B_i\cup E(H_i) \backslash \{h\}$ will color $V(IG(H_i))$.
    With this, it can be concluded that $B_i\cup E(H_i) \backslash \{h\}$ is indeed a zero forcing set for $IG(H_i)$.
    \item[Case $B_i=\emptyset$.] This case is only possible if there is at least one singleton hyperedge $h\in E(H_i)$. Then consider the vertex $v\in h$. Since $H_i$ is connected and there is at least one non-singleton hyperedge, there is a non-singleton hyperedge $h'\in E(H_i)$ such that $v\in h'$. 
    Then consider the set $E(H_i)\backslash\{h'\}$ (as $B_i=\emptyset$). In the first round of the zero forcing process on $IG(H_i)$, the force $h\rightarrow v$ will occur and in the second round the force $v\rightarrow h'$ will occur.
    Therefore, at least $E(H_i)$ will be colored after the second round of the zero forcing process and hence $E(H_i)\backslash \{h'\}$ is a zero forcing set.
    
    From this, the following inequality is obtained for $1\leq i \leq k$:
    
    \[Z(IG(H_i)) \leq b_L(H_i) + |E(H_i)| - 1.\]
    
    Additionally, inequality~\eqref{eq:hypergraph inequality} can be applied to each of the connected component $H_i$ for $k < i \leq l$ of $H$. Then, summing over all the connected components we obtain
    \begin{align*}
        Z(IG(H)) &= \sum_{i=1}^{k} Z(IG(H_i)) + \sum_{i=k+1}^l Z(IG(H_i)) \\
        &\leq \sum_{i=1}^{k} (b_L(H_i)+|E(H_i)|-1) + \sum_{i=k+1}^l (b_L(H_i)+|E(H_i)|) \\
        &= b_L(H) + |E(H)| - k.
 \end{align*} 
\end{description}
This completes the proof.
\end{proof}

\end{document}
